\backhead{OBJECTIVE TYPE QUESTIONS}{khred}
\begin{exlist}
\item The reference evapotranspiration of FAO-56 is the demand of\opt{bare soil}{a well-watered grass}{open water}{a forest}
\item The daily damping depth of a moist soil is of the order of\opt{1 mm}{1 cm}{10 cm}{1 m}
\item The degree-day factor of snowmelt is about\opt{0.1}{3}{30}{300} mm\,$^\circ$C$^{-1}$\,d$^{-1}$
\item Capillary rise from a water table falls with depth\opt{linearly}{exponentially}{not at all}{as $\sqrt L$}
\item The factor 0.408 in the Penman--Monteith equation converts\opt{kPa to mm}{MJ\,m$^{-2}$ to mm}{K to $^\circ$C}{m\,s$^{-1}$ to km\,d$^{-1}$}
\end{exlist}
\objkey{1.\ (b),\quad 2.\ (c),\quad 3.\ (b),\quad 4.\ (b),\quad 5.\ (b)}
\dykband{EXERCISES IN SWAFL}
{\small Open SWAFL (desktop or online) with the benchmark project. The chapter of the \href{https://github.com/halilunaziru73-creator/Coupled-Water-Fluxes-Pore-Scale-Dynamics-Precision-Irrigation-Mulching-Feedback-Loops/raw/main/Software/v1.1.0/SWAFL_Desktop/docs/manual_v1.1.0/SWAFL_Manual_v1.1.0.pdf}{SWAFL manual} that explains each task is given in brackets; record your values in a table and compare them with what is expected (\emph{Check}).}\par\vspace{3pt}
\begin{exlist}
\item {\color{khblue}[Manual, Chapter 25]} Open the Daisy process map and click every box of the weather, water and heat part. For each, write the equation it uses.\par\noindent{\small\itshape Check: each box opens its description and equation.}
\item {\color{khblue}[Manual, Chapter 19]} Use \emph{The weather of a place} to download the weather of your town and check the radiation and temperature.\par\noindent{\small\itshape Check: the record covers the twenty-year grid.}
\item {\color{khblue}[Manual, Chapter 30]} Open the plot of the soil temperature of a run and read the annual amplitude at two depths.\par\noindent{\small\itshape Check: the amplitude falls with depth.}
\item {\color{khblue}[Manual, Chapter 30]} Read the water balance of a run in the raw output and close it for one year.\par\noindent{\small\itshape Check: the residual is small.}
\end{exlist}
